\scanpage{539}
\mspage{450}
% archive sheet number: 538

Posons alors

\[p_i=\alpha_i\beta_i:\unsure{x_i|z_i}\mathrel{\unsure{\relbar|\to}}x_i,\]

il est clair qu'on a

\[\left\{\begin{aligned}&p_i^2=p_i\quad(p_i\text{ un projecteur})\\ &\alpha_i:z\longrightarrow x_i\quad\text{induit un isom.}\\ &\hspace{2em}z\xrightarrow{\sim}x_i^\natural\overset{\mathrm{déf.}}{=}\operatorname{Im}p_i=\operatorname{Ker}(\mathrm{id}_{x_{\unsure{j|i}}},p_i).\end{aligned}\right.\]

L'ens \inserted{ordonné} cofinal cherché $J$ est $i_0\backslash I_{\deleted{\unsure{\geq i_0}}}$, il est filtrant, et est \unsure{grd|gd|pt|petit} devant $\pi_0$ \unsure{si} $I$ l'est. \deleted{Il faut vérifier les conditions $1^\circ$ \unsure{;|à} $4^\circ$.}

À vrai dire, la réduction au cas $I$ ordonné est inutile, on peut procéder directement $J=i_0\backslash I$, Il est immédiat que $\deleted{\unsure{I}}J$ est filtrante ($I$ l'étant) et $i_0\backslash I\mathrel{\relbar}I$ cofinal, enfin que si $I$ est grd devant $\pi_0$\inserted{\unsure{un cardinal infini|un ordre infini}}, il en est de $=$ de $i_0\backslash I$. \deleted{\unsure{Condition|Considérons} \unsure{$1^\circ$|$1^\circ$ :}}

Exprimons que $\beta\alpha=\mathrm{id}_{\underline{x}}$ ; or on voit que $\beta\alpha$ est défini par l'endomorphisme $(p_j)_{j\in J}$ du système inductif $(x_j)_{j\in J}$ --- le fait que ce soit bien un endom. de ce système, i.e. qu'on ait la relation $1^\circ)$ de 4.15.7.1.

\[p_{j'}x_\lambda=x_\lambda p_{\unsure{\lambda|j}}\qquad\text{si }\lambda:j\mathrel{\relbar}j'\text{ ds }J,\]

est immédiat, ça s'écrit

\[\alpha_{j'}\underbrace{\beta_{j'}x_\lambda}_{\beta_j}=\underbrace{x_\lambda\alpha_j}_{\alpha_{j'}}\beta_j\qquad\text{or les deux mbs sont égaux à }\alpha_{j'}\beta_j.\]
